\section{Validation of Design Science Control Principles and Heuristics}
\label{sec:dsr-validation}

A fundamental requirement of the Design Science Research (DSR) methodology is ensuring that the theoretical constructs defined in the Knowledge Cycle are tangibly materialized in the instantiated artifact. The experimental architecture and the automated orchestration scripts successfully validated the Control Principles (CP) and Contingency Heuristics (ConH) proposed in Chapter 5.

\textbf{Materialization of Control Principles}

The Control Principles dictate the structural axioms required to ensure the telemetry pipeline behaves as a reliable feedback loop.

\begin{itemize}
    \item \textbf{CP1 (Principle of Causal Linkage):} To prevent causal ambiguity during latent failures, signals must preserve temporal and topological context. This principle was technically validated through the Infrastructure as Code (IaC) configuration (Appendix B). The Ansible deployment natively injects \texttt{derivedFields} into the Loki log aggregation backend, mapping the regex-extracted \texttt{TraceID} directly to Grafana Tempo. Simultaneously, the Prometheus configuration enforces \texttt{send\_exemplars: true}, structurally linking aggregated metric spikes to specific distributed trace spans. This guarantees that diagnostic transitions between metrics, logs, and traces are deterministic, eliminating the need for manual, time-bound cross-referencing.
    
    \item \textbf{CP2 (Principle of Unified State Representation):} To overcome the fragmentation of agent deployment, telemetry signals must conform to a standardized schema. This was validated by the \texttt{panoptes-collector.yml} Custom Resource. Instead of deploying disparate agents for each pillar, the OpenTelemetry Collector was configured with a single \texttt{otlp} receiver handling metrics, logs, and traces concurrently over gRPC port 4317. This unified representation is what allowed the automated RCA engine (Scenario 4) to parse the JSON call-graph deterministically without relying on vendor-specific APIs.
\end{itemize}

\textbf{Materialization of Contingency Heuristics}

The Contingency Heuristics define the adaptive rules required to maintain internal validity and operational stability during the experiment.

\begin{itemize}
    \item \textbf{ConH1 (Heuristic of Declarative Control):} This heuristic mandates that observability systems must adapt to ephemeral workloads without manual intervention. This was realized through the deployment of the \texttt{ServiceMonitor} Custom Resource (Appendix B). By dynamically discovering the Astronomy Shop microservices via label selectors rather than static IP configurations, the Prometheus Operator ensured that all newly scheduled pods were immediately scraped, demonstrating architectural resilience to dynamic scaling.
    
    \item \textbf{ConH2 (Heuristic of Configuration Drift):} To guarantee that the measured diagnostic latency ($MTTDiag$) was solely attributable to the architecture rather than operator variance, state changes were treated as strict external disturbances. The orchestration scripts (\texttt{run-scenario4-optimized-trials.py}) enforced this by clearing residual traces and executing a programmatic cluster synchronization (\texttt{kubectl rollout status}) prior to every Chaos Mesh injection. The completely automated nature of the $N=30$ iterations ensured zero configuration drift across the experimental bounds.
    
    \item \textbf{ConH3 (Heuristic of Adaptive Feedback):} To prevent the observability stack from degrading into a "noisy neighbor" under heavy monitoring loads, this heuristic restricts the overhead envelope. This was effectively validated by configuring the \texttt{batch} processor within the OpenTelemetry Collector pipeline. By buffering telemetry data into discrete chunks prior to transmission, the architecture successfully minimized CPU contention, allowing Hypothesis $H_{1,3}$ to be accepted with the total PANOPTES CPU overhead consuming only $16.75\%$ of the cluster capacity.
\end{itemize}